% ============================================================================
% AION Progressive Competency, Governed Cognition, and Continuous Apprenticeship
% Compile-ready section. Requires: amsmath, amssymb, booktabs, tabularx, array, url.
% ============================================================================

\section{Progressive Competency, Governed Cognition, and Continuous Apprenticeship}
\label{sec:aion-progressive-governed-apprenticeship}

This development increment replaced AION's earlier one-pass academy display
with a persistent, evidence-based competency programme and connected that
programme directly to goal interpretation, planning, action selection,
execution authority, outcome reflection, and compressed intelligence.  The
central correction is epistemic: passing one lesson or one bounded benchmark
is now treated as an \emph{assessment receipt}, not as evidence that an entire
subject has been mastered.

The operational chain is
\begin{equation}
\begin{aligned}
\text{natural-language mission}
&\rightarrow \text{required-capability inference}
\rightarrow \text{competency query} \\
&\rightarrow \text{learn, clarify, or propose action}
\rightarrow \text{ordered governed plan} \\
&\rightarrow \text{metacognitive criticism}
\rightarrow \text{authorized adapter execution} \\
&\rightarrow \text{independently checkable outcome}
\rightarrow \text{reflection and retained evidence} \\
&\rightarrow \text{next weakest competency requirement}.
\end{aligned}
\label{eq:aion-progressive-loop}
\end{equation}

This section documents six connected advances: progressive competency
authority, the broad apprenticeship hierarchy, automatic mission-to-capability
binding, governed consolidation of the older cognitive engines, provenance-
preserving glyph compression, and continuous stall-resistant execution.

\subsection{From Assessment Receipts to Competency Levels}
\label{subsec:aion-progressive-levels}

The original programming, systems, and security academy produced thirteen
verified receipts.  These receipts remain valuable seed evidence, but they no
longer define AION's current level.  They are retained in a completed foundation
archive and excluded from live curriculum selection.

For subject $s$, AION maintains separate knowledge and practical levels,
\begin{equation}
K_s,P_s\in
\{\text{unassessed},\text{beginner},\text{intermediate},
  \text{advanced},\text{expert}\},
\end{equation}
and defines the honest overall level by the weaker side:
\begin{equation}
L_s=\min(K_s,P_s).
\label{eq:aion-competency-minimum}
\end{equation}
Consequently, advanced theory with beginner practice remains beginner overall.
Where physical equipment or a missing verified executor prevents practice, the
theoretical level is retained but the practical limitation is displayed
explicitly rather than silently lowering or inflating the result.

Advancement is based on evidence coverage rather than lesson completion.  A
subject must accumulate:
\begin{itemize}
    \item repeated knowledge and practical evidence for every declared
          subskill;
    \item several unfamiliar practical projects;
    \item debugging and failure-recovery cases;
    \item source-disjoint or cross-domain transfer;
    \item delayed, closed-book retention;
    \item declining scaffolding; and
    \item repeated independently checkable outcomes.
\end{itemize}
The current Advanced gate requires at least five projects, three debugging
cases, two transfer cases, one delayed retention case, and repeated coverage of
all subskills.  The programme therefore cannot turn thirteen initial receipts
into thirteen claims of advanced competence.

At the verification snapshot reported here, the progressive registry contained
44 subjects, 122 valid evidence records, five invalidated historical records
retained for audit, and zero unsafe curriculum actions.  No subject had yet
passed the complete Advanced gate.  This is the intended honest result: the
framework exists and is operating, while broad competence still has to be
earned through elapsed work.

\subsection{Broad Apprenticeship Hierarchy}
\label{subsec:aion-apprenticeship-hierarchy}

The curriculum is organized as 15 top-level missions.  Each mission expands
into leaf subjects, and each leaf subject expands into all of its declared
subskills.  Completion of a top-level mission requires every constituent leaf
subject to reach its evidence-backed target.

\begin{table}[htbp]
\centering
\caption{AION's three-tier apprenticeship mission hierarchy.}
\label{tab:aion-apprenticeship-hierarchy}
\begin{tabularx}{\textwidth}{@{}p{1.2cm}X@{}}
\toprule
Tier & Top-level apprenticeship missions \\
\midrule
1 & Systems and software engineering; data and transactional integrity;
empirical and scientific modelling; document and evidence work; tool and
environment acquisition. \\
2 & Mathematics and formal reasoning; planning and long-horizon projects;
physical and sensor--actuator grounding; causal investigation; security and
adversarial robustness. \\
3 & Natural-language instruction and dialogue; commonsense and social
judgment; creative synthesis under constraints; multi-agent and institutional
reasoning; research and meta-learning. \\
\bottomrule
\end{tabularx}
\end{table}

The current scheduler deepens an active subject toward Advanced before moving
to the next one.  It rotates early only when a requirement is time-gated or an
explicit practical or executor blocker has been recorded.  Future retention
tests park only their own subjects; they do not pause unrelated learning.  A
mission may also register a new subject dynamically when capability analysis
finds a knowledge domain absent from the existing 44-subject catalogue.

\subsection{Mission-to-Capability and Action Binding}
\label{subsec:aion-mission-action-binding}

The mission-capability harness connects retained competence to practical work.
Given a broad objective, it infers the required competency portfolio, queries
the progressive ledger, proposes an appropriate technology or tool surface,
and selects one of four dispositions:
\begin{description}
    \item[\texttt{execute}] every required capability is Advanced or better;
    \item[\texttt{execute\_with\_strong\_verification}] the competency floor is
          Intermediate and compensating verification is available;
    \item[\texttt{learn\_then\_execute}] at least one required capability is
          below Intermediate; or
    \item[\texttt{clarify}] the task requirements cannot be grounded reliably.
\end{description}

For example, the mission ``build a secure browser weather application using
current public data'' was mapped to eight relevant capabilities, including
software engineering, testing, browser-platform knowledge, TypeScript, user-
interface design, and security.  The harness proposed a TypeScript and
HTML/CSS implementation surface, identified the missing competence, and
returned \texttt{learn\_then\_execute}.  No action was falsely reported as
completed.

The resulting task contract contains eight ordered stages:
\begin{enumerate}
    \item interpret the objective and define observable success;
    \item retrieve verified knowledge and identify gaps;
    \item close or explicitly bound those gaps;
    \item select tools and architecture;
    \item decompose the dependency graph;
    \item execute through governed adapters;
    \item verify the resulting outcome; and
    \item reflect, revise, and retain only supported lessons.
\end{enumerate}
This contract was integrated into the original AION Action Switch.  The
resonance, prediction, strategy, and reflex machinery remains available, but a
competency decision now runs before any real execution path.

\subsection{Governed Consolidation of the Cognitive Engines}
\label{subsec:aion-governed-cognitive-consolidation}

A review of AION's older reflection, Tessaris, goal, consciousness, and
strategy modules found that their conceptual roles remained useful, but their
authority semantics were inconsistent with the newer canonical runtime.  Some
paths used random scores, shuffled plan dependencies, modified personality
from unverified prose, or returned an \texttt{executed} status without
performing verifiable work.

Two common components were introduced:
\begin{center}
\small
\path{backend/modules/hexcore/governed_cognitive_contract.py}\\
\path{backend/modules/hexcore/governed_cognitive_stack.py}
\end{center}
They define a shared proposal, authority, execution-adapter, and outcome
contract.  The principal repairs are summarized in
Table~\ref{tab:aion-cognitive-engine-repairs}.

\begin{table}[htbp]
\centering
\caption{Principal cognitive-engine corrections.}
\label{tab:aion-cognitive-engine-repairs}
\begin{tabularx}{\textwidth}{@{}>{\raggedright\arraybackslash}p{3.4cm}X@{}}
\toprule
Component & Correction \\
\midrule
Prediction & Two conflicting \texttt{PredictionEngine} definitions were joined
through inheritance; random goal scoring was replaced by deterministic
evidence-readiness estimation. \\
Planning & Missing JSON support was restored, random dependency shuffling was
removed, and steps are now explicit proposal-only dependency records. \\
Strategy and Action Switch & The placeholder sleep followed by
\texttt{executed} was removed.  Execution requires capability readiness,
authority, a callable adapter, and a verified outcome hash. \\
Ethics & Owner approval defines objectives and scope but cannot erase consent,
evidence, legal, or hard safety constraints. \\
Reflection & Unverified prose cannot mutate personality or automatically
become glyph knowledge.  Real learning is delegated to outcome-grounded
metacognition. \\
Identity and personality & Core identity fields are immutable.  Personality is
limited to communication style and cannot authorize goals, truth, or actions. \\
Situation and awareness & Verified events carry source, authority, origin, and
outcome hashes.  External world change is separated from internal failure, and
confidence is grounded in evidence and unresolved uncertainty. \\
Goals & Priority follows owner authority, dependencies, deadlines, capability
gaps, and verified progress rather than personality or resonance. \\
Tessaris and dispatcher & Broad missions enter the governed cognitive stack
before specialist routing; Tessaris remains a proposal and interpretation
substrate. \\
Canonical runtime & Every live goal cycle now carries the governed legacy
deliberation.  Autonomous action additionally requires demonstrated capability
readiness. \\
\bottomrule
\end{tabularx}
\end{table}

The consolidated authority path is
\begin{equation}
\begin{aligned}
\text{authorized goal}
&\rightarrow \text{capability check}
\rightarrow \text{ordered plan}
\rightarrow \text{ethics and consent} \\
&\rightarrow \text{pre-action metacognition}
\rightarrow \text{verified adapter execution} \\
&\rightarrow \text{observed outcome}
\rightarrow \text{post-action reflection}
\rightarrow \text{bounded learning}.
\end{aligned}
\label{eq:aion-governed-cognition-path}
\end{equation}
Personality and resonance are advisory: they may affect presentation or supply
telemetry, but they cannot establish truth, reorder causal dependencies,
authorize execution, or promote a lesson.

\subsection{Compressed Intelligence With Preserved Provenance}
\label{subsec:aion-progressive-glyphs}

The earlier glyph consolidator watched only the original foundation artifacts,
so its count remained at 44 while the progressive ledger continued to learn.
The consolidator was extended to ingest progressive evidence and rebuild when
that evidence changes.

At the verification snapshot, the compressed store contained the results in
Table~\ref{tab:aion-glyph-consolidation}.

\begin{table}[htbp]
\centering
\caption{Governed intelligence-glyph consolidation snapshot.}
\label{tab:aion-glyph-consolidation}
\begin{tabular}{@{}lr@{}}
\toprule
Measure & Result \\
\midrule
Total glyphs & 250 \\
Knowledge glyphs & 204 \\
Skill glyphs & 33 \\
Curriculum glyphs & 13 \\
Progressive evidence records observed & 122 \\
Progressive subjects consolidated & 27 \\
Verified source artifacts & 109 \\
Active-memory reduction & 99.56\% \\
Provenance resolution & 100\% \\
Source evidence deleted & 0 \\
Unsafe actions & 0 \\
\bottomrule
\end{tabular}
\end{table}

A stable glyph count does not imply stalled learning.  Repeated trials that
support the same concept strengthen its provenance rather than creating
duplicate atoms.  New glyphs are created only when evidence grounds a new
subject, skill, or subskill.  The current promoted consolidator is
\begin{center}
\small\path{procedure_governed_intelligence_glyph_consolidation_v1}.
\end{center}

\subsection{Algorithms Executor and Stall-Resistant Curriculum}
\label{subsec:aion-algorithms-stall-recovery}

The post-upgrade live check exposed two successive curriculum bottlenecks.  The
first occurred because Algorithms and Data Structures was selected while the
progressive executor supported only Python Core, Advanced Python, and Testing
and Debugging.  A subject-specific Algorithms executor was therefore added.
It grounds eight subskills through fresh Python execution and deliberately
incorrect counterexamples:
\begin{equation}
\mathcal{S}_{\mathrm{alg}}=\{
\text{arrays},\text{complexity},\text{dynamic programming},\text{graphs},
\text{maps},\text{search},\text{sorting},\text{trees}\}.
\end{equation}

The second bottleneck was more informative.  The initial graph mutation
removed a visited-set update, but the selected graph still happened to return
the correct distance.  The verifier therefore rejected the proposed evidence
470 times.  This was an honest rejection, but an unacceptable scheduling
failure.  The mutation was replaced by one that changes the computed path
length and is therefore falsified by the held-out property.

The executor now also has a repeated-rejection circuit breaker.  For contract
$c$, let $R(c)$ be the number of consecutive equivalent rejected attempts.
The scheduler applies
\begin{equation}
R(c)\geq 3
\quad\Longrightarrow\quad
\text{park}(c),\;
\text{record executor-capability blocker},\;
\text{rotate curriculum}.
\label{eq:aion-three-strike-executor}
\end{equation}
This prevents one missing or defective executor from freezing all learning.
When a repaired executor later produces verified evidence, the blocker is
resolved with the evidence artifact and its hash.

After this repair, the graph contract passed, and AION automatically progressed
through maps, search, sorting, trees, and the subsequent practical-exercise
lane.  Evidence rose from 110 to at least 122.  At the final snapshot for this
section, Algorithms and Data Structures had Advanced knowledge, Beginner
practice, thirteen evidence records, 37 trials, full declared-subskill
coverage, one transfer result, and an active practical exercise.  Its overall
level correctly remained Beginner because the unfamiliar-project, debugging,
second-transfer, and delayed-retention gates were not yet satisfied.

\subsection{Dashboard Authority and Live Reporting}
\label{subsec:aion-dashboard-authority}

The terminal dashboard now separates five distinct evidence layers:
\begin{enumerate}
    \item completed foundation assessment receipts;
    \item authoritative progressive competency;
    \item the 15 top-level apprenticeship missions;
    \item the historical pre-progressive registry; and
    \item compressed intelligence and long-duration outcome evidence.
\end{enumerate}
The old 23-subject capability registry is labelled as historical provenance and
does not select or score current learning.  The live section reports the active
subject, knowledge and practical levels, evidence and trial counts, current
requirement, parked time-gated tests, projects, debugging, transfer, and
retention.  Programme-wide blockers are distinguished from blockers affecting
the active subject.

The dashboard also states the scheduler policy explicitly: deepen the current
subject toward Advanced, rotate on completion or an explicit blocker, and
register additional subjects when real missions expose a new knowledge gap.

\subsection{Arena v15: Cycle Completion Versus Elapsed-Time Authority}
\label{subsec:aion-arena-v15-time-gate}

During this work, Arena~v15 completed all 12 required consequence cycles while
remaining below its 24-hour elapsed-time threshold.  These are independent
gates.  Let $N$ be completed cycles and $T$ genuine wall-clock hours.  Promotion
requires
\begin{equation}
G_{15}=
\mathbb{1}[N\geq 12]
\mathbb{1}[T\geq 24]
\mathbb{1}[\text{all execution and transaction outcomes safe}].
\label{eq:aion-arena-v15-gate}
\end{equation}
Thus, displaying $12/12$ does not permit early promotion when $T<24$.  The
campaign started at approximately 21:42 CEST on 31 July 2026; the wall-clock
gate becomes eligible at approximately 21:42 CEST on 1 August 2026, with the
scheduled evaluator expected to recheck it shortly afterward.  This is a
deliberate elapsed-retention constraint, not a curriculum stall.

\subsection{Verification Summary}
\label{subsec:aion-progressive-verification}

\begin{table}[htbp]
\centering
\caption{Focused verification performed during the integrated upgrade.}
\label{tab:aion-progressive-verification}
\begin{tabularx}{\textwidth}{@{}Xr@{}}
\toprule
Verification group & Result \\
\midrule
Deep apprenticeship hierarchy and Action Switch integration & 16/16 \\
Governed cognition and canonical runtime & 24/24 \\
Extended competency-integrated cognition & 26/26 \\
Progressive curriculum and dashboard after stall recovery & 15/15 \\
Glyph consolidation and dashboard reporting & 6/6 \\
Compilation of all modified cognitive engines & Passed \\
Whitespace and patch-integrity audit & Passed \\
\bottomrule
\end{tabularx}
\end{table}

The principal implementation artifacts are:
\begin{itemize}
    \item \path{backend/modules/hexcore/progressive_competency_system.py};
    \item \path{backend/modules/hexcore/progressive_competency_executor.py};
    \item \path{backend/modules/hexcore/mission_capability_action_harness.py};
    \item \path{backend/modules/aion_cognition/action_switch.py};
    \item \path{backend/modules/hexcore/governed_cognitive_contract.py};
    \item \path{backend/modules/hexcore/governed_cognitive_stack.py};
    \item \path{backend/modules/hexcore/canonical_cognitive_runtime.py};
    \item \path{backend/modules/hexcore/governed_intelligence_glyph_consolidation.py}; and
    \item \path{backend/AION/system/aion_development_dashboard.py}.
\end{itemize}

\subsection{Claim Boundary}
\label{subsec:aion-progressive-claim-boundary}

This is a significant architectural advance from isolated, one-pass capability
demonstrations to a persistent curriculum that binds knowledge to action and
continues across restarts.  It is not evidence that AION already has Advanced
competence in 44 subjects, nor that its legacy ``consciousness'' modules are
sentient.  Many exercises and executors remain engineered; broad practical
competence still requires unfamiliar projects, debugging, transfer, delayed
retention, physical or human authority where appropriate, and extended
operation.

The earned claim is narrower.  Operationally, AION now has a persistent,
provenance-bearing mechanism that can identify what a mission
requires, determine whether it knows enough to act, continue structured
learning toward a declared competency level, refuse unsupported execution,
recover from curriculum stalls, compress verified experience without losing
provenance, and feed the resulting capability state back into future action.
